%%%%%%%%%%%%%%%%%%%%%%%%%%%%%%%%%%%%%%%%%%%%%%%%%%%%%%%%%%%%%%%%%%%%%%
%  月度报告分享稿（控制算法方向）—— 讲稿版
%
%  与报告模板不同：这是一份「照读用」的讲稿，所以版式偏易读
%    · 大行距，方便边看边讲
%    · 有「提示框」标出重点页、时间、注意事项
%    · 不追求正式排版，追求看得清
%
%  编译：xelatex 本文件（跑两遍）
%  Overleaf：Menu -> Compiler -> XeLaTeX
%%%%%%%%%%%%%%%%%%%%%%%%%%%%%%%%%%%%%%%%%%%%%%%%%%%%%%%%%%%%%%%%%%%%%%
\documentclass[11pt,a4paper]{ctexart}

\usepackage[margin=2.0cm]{geometry}
\usepackage[table]{xcolor}
\usepackage{booktabs}
\usepackage{tabularx}
\usepackage{array}
\usepackage{ragged2e}
\usepackage{enumitem}
\usepackage{titlesec}
\usepackage{fancyhdr}
\usepackage{graphicx}
\usepackage{amsmath}
\usepackage{amssymb}
\usepackage[hidelinks]{hyperref}
\usepackage{tcolorbox}
\tcbuselibrary{skins,breakable}

% ---------- 颜色 ----------
\definecolor{cBlue}{HTML}{0B4C8C}
\definecolor{cBlueII}{HTML}{1868A8}
\definecolor{cGrey}{HTML}{5A5A5A}
\definecolor{cLight}{HTML}{F2F7FB}
\definecolor{cWarnBg}{HTML}{FFF8E1}
\definecolor{cWarnLine}{HTML}{E6A23C}
\definecolor{cKeyBg}{HTML}{EAF3FA}
\definecolor{cRed}{HTML}{C62828}
\definecolor{cGreen}{HTML}{2E7D32}

% ---------- 版式 ----------
\linespread{1.45}
\setlength{\parindent}{0pt}
\setlength{\parskip}{6pt}
\newcolumntype{Y}{>{\RaggedRight\arraybackslash}X}
\renewcommand{\arraystretch}{1.3}

\titleformat{\section}{\Large\bfseries\color{cBlue}}{\thesection}{0.6em}{}
\titlespacing{\section}{0pt}{14pt}{6pt}
\titleformat{\subsection}{\large\bfseries\color{cBlueII}}{\thesubsection}{0.6em}{}
\titlespacing{\subsection}{0pt}{10pt}{4pt}

\pagestyle{fancy}
\fancyhf{}
\fancyhead[L]{\small\color{cGrey}月度报告分享稿 · 控制算法}
\fancyhead[R]{\small\color{cGrey}张昊洋}
\fancyfoot[C]{\small\color{cGrey}第 \thepage\ 页}
\renewcommand{\headrulewidth}{0.4pt}

\setlist[itemize]{leftmargin=1.6em,itemsep=3pt,topsep=4pt}
\setlist[enumerate]{leftmargin=2.0em,itemsep=3pt,topsep=4pt}

% ---------- 提示框 ----------
\newtcolorbox{tishi}[1][]{%
  breakable, enhanced, colback=cWarnBg, colframe=cWarnLine,
  boxrule=0.9pt, arc=2pt, left=7pt, right=7pt, top=5pt, bottom=5pt,
  fontupper=\small, #1}
\newtcolorbox{yaodian}[1][]{%
  breakable, enhanced, colback=cKeyBg, colframe=cBlueII,
  boxrule=0.9pt, arc=2pt, left=7pt, right=7pt, top=5pt, bottom=5pt,
  fontupper=\small, #1}

% 「照读」标记：实际要说出口的话
\newcommand{\shuo}[1]{\par{\color{cBlue}\bfseries 说：}{\itshape #1}\par}
% 提示动作
\newcommand{\dongzuo}[1]{\par{\color{cWarnLine}\bfseries 【#1】}\par}

\begin{document}

\begin{center}
{\LARGE\bfseries\color{cBlue} 月度报告分享稿}\\[6pt]
{\large 控制算法方向}\\[10pt]
{\small\color{cGrey}
汇报人：张昊洋\hspace{1em}｜\hspace{1em}方向：流量动态调控\hspace{1em}｜\hspace{1em}两份报告：阶段一 8 月（6 页）＋ 阶段二 9 月（7 页）}\\[4pt]
{\small\color{cGrey}建议时长 20--25 分钟（阶段一 8 分钟，阶段二 14 分钟，留 3 分钟提问）}
\end{center}

\vspace{4pt}

\begin{tishi}
\textbf{怎么用这份稿子}\\[3pt]
· \textbf{蓝色「说：」后面的斜体字}是可以直接讲出口的话，不用背，照读就顺。\\
· \textbf{橙色【】}是动作提示，提醒你切页面、指图。\\
· \textbf{两个[重点]标记的部分是重点}，一定要留够时间；时间不够先砍阶段一的「核心概念」。\\
· 附录 A 是预计被问到的问题，附录 B 是时间分配和三个提醒。
\end{tishi}

\section{开场}

\dongzuo{30 秒，不要拖}

\shuo{各位老师、同学好，我汇报控制算法这一部分的两次月度总结。}

\shuo{先说一下这两份报告的分工。阶段一是初步阅读，主线是「我读懂了什么」，
重点是文献判断和一个复现实验；阶段二是研究推进，主线是「我做了什么、发现了什么」，
重点是仿真结果和组长布置的控制实验方案。}

\shuo{两份报告是递进关系——阶段一读文献得出的判断，直接决定了阶段二做什么。}

\section{阶段一 · 初步阅读}

\subsection{阅读范围}

\dongzuo{1.5 分钟，指文献清单表}

\shuo{我这个月读了 5 篇文献，选读的线索是两条。}

\shuo{第一条线索是判断「要不要上智能控制」。先读 Mnih 的 Nature 文章和 PPO 原文，
把深度强化学习的方法框架搞清楚；然后重点读 Vignon 2023 那篇综述，
看 DRL 在流动控制领域实际做到了什么程度。}

\shuo{第二条线索是「物理模型和算法怎么接起来」。读陶飞院士的数字孪生和 Brunton 的
《数据驱动科学与工程》，明确「物理模型打底、数据修正」的路径。}

\shuo{我没有读纯控制理论教材，因为本科阶段已经有基础，现阶段更缺的是判断力，不是推导能力。}

\subsection{核心概念}

\dongzuo{2 分钟，讲第一个概念时切到图 1}

\begin{yaodian}
\textbf{第一个概念，也是我这个月最大的收获：}\\
\textbf{执行器到被控量的映射，才是控制问题的真正起点。}
\end{yaodian}

\dongzuo{指着图 1}

\shuo{这张图是我们系统的变量结构。左边是\textbf{我们能调的}——泵转速，未来还要加个水阀；
中间是流速，右边是芯片温度。}

\shuo{问题在于：从泵转速到芯片温度，中间隔着两段物理过程——「泵到流速」，
和「流速到散热片的换热」。}

\shuo{这两段我们都不清楚。不清楚意味着什么？意味着我不知道「改变一个单位的泵转速，
温度会变多少」——那就无法确定 PID 的比例系数该取多大。}

\shuo{所以我们把系统整理成了两个待辨识的函数，就是图上标的 f 和 g。
这也是组长后来说要做控制实验的原因。}

\shuo{第二个概念是液冷系统的大热惯性，它决定了控制带宽的上限。
惯性来自三处：散热片的热容、冷却液自己的储热、还有流体从泵流到芯片的输运延迟。}

\shuo{惯性大的意思是——即使我立刻把泵开到最大，温度也不会立刻降下来。
所以纯反馈控制天然带滞后，一味提高增益会导致超调和振荡。}

\shuo{第三个概念是增量式 PID 和积分饱和。增量式 PID 输出的是控制量的增量而不是绝对值，
工程上有两个好处：手动切自动时冲击小，也不容易因为积分累积而失控。
但积分饱和要单独处理——输出到了行程上限后，误差还在累积积分项，
反向调节时就会严重滞后。}

\subsection{复现}

\dongzuo{2.5 分钟，切到图 2}

\shuo{阶段一的硬要求是「必须有一个复现」——只列文献不算完成。}

\shuo{我选的复现对象是二阶线性系统的阶跃响应。选它是因为它是控制理论里最基础、
最容易核验的模型——如果连这个都复现不准，后面复杂模型的可信度就没法谈了。}

\dongzuo{指图 2}

\shuo{这是六个阻尼比下的响应曲线。可以看到：阻尼比 0.3 的时候超调 37\%，
0.707 的时候只有 4.33\%。0.707 是超调和快速性的常用折中，这个结论后面调 PID 时会用到。}

\shuo{理论值和仿真量取值完全吻合，误差在采样精度以内。}

\shuo{但复现里发现了一个有意思的事情。}

\dongzuo{切到图 3，指右图}

\shuo{调节时间的理论公式是 $t_s = 4/(\zeta\omega_n)$。我用仿真曲线直接量取了实际调节时间，
发现在阻尼比小于 0.7 时，这个公式偏保守——比如阻尼比 0.3 时，公式给 13.33 秒，
实际量取约 11.8 秒，差了 11\%。}

\begin{yaodian}
\textbf{这个发现看起来很小，但直接影响了我的做法：}\\
后续评价控制器性能时，我改用\textbf{仿真量取值}，不再直接套公式。
\end{yaodian}

\subsection{收获与疑问}

\dongzuo{2 分钟}

\begin{yaodian}
\textbf{这个月我认知上最大的变化是：}\\
\textbf{控制效果的瓶颈，往往不在算法，而在我们对控制通道的了解程度。}
\end{yaodian}

\shuo{我原来以为，控制就是「选一个更好的算法」。读完 Vignon 的综述以后发现——
很多 DRL 控制研究之所以效果好看，是因为有廉价的仿真环境可以无限交互。
而我们一旦涉及真实热惯性、流固耦合、传感器噪声，仿真和实物的差距就会成为主要误差源。}

\shuo{所以我主动把重点从「学更复杂的算法」转成了「先把控制通道走通」。
这个判断和组长后来布置的控制实验是一致的方向。}

\shuo{还有三个没搞清楚的问题，我列在报告里：}

\begin{enumerate}
  \item 组长提出未来要加\textbf{水阀}，那就会出现「泵 + 阀」两个执行器。
        \textbf{它们会不会相互耦合？}比如关小阀门后，泵的工作点会不会移动？
        如果耦合明显，就不能当独立变量处理。
  \item 液冷回路的热惯性\textbf{到底有多大}？仿真里用的时间常数是我\textbf{假定}的，
        没有实测依据。如果实测远大于假定值，整套参数都要重来。
  \item 传感器只有离散测点，但要稳的是芯片最高温。
        \textbf{用测点代表最高温，误差有多大？}这决定了控制目标怎么定义。
\end{enumerate}

\section{阶段二 · 研究推进}

\dongzuo{切换到阶段二报告}

\shuo{阶段二的主线变成「研究推进」。我按任务、阅读、进展、问题、结论、计划六个部分讲。}

\subsection{任务与完成情况}

\dongzuo{1.5 分钟，指完成情况表}

\shuo{月初我定了三个目标，完成情况是这样：控制实验方案整理、控制回路建模、
基础 PID 仿真验证、智能算法可行性判断，这四项都已完成。}

\begin{tishi}
\textbf{要主动说明的一点：}\\
「执行器/散热片实测辨识」这一项是\textbf{未开始}的，因为现在还没有实验条件。
\textbf{这一条我没做，不能算完成。}\\[3pt]
—— 主动说比被问出来好。
\end{tishi}

\subsection{把组长的要求整理成模型}

\dongzuo{2 分钟，切到图 1}

\shuo{这是本月最核心的一项工作——把组长提出的控制实验，整理成可执行的方案。}

\shuo{组长说：热流密度是我们无法控制的变量；入口水温和流速是自变量；
散热片是个复杂的函数；芯片温度是因变量，我们要让它稳定。}

\shuo{我把这个整理成三条链：}

\dongzuo{指图}

\shuo{第一，变量角色的划分——热流密度是不可控扰动，芯片负载决定，我们改不了；
入口水温、泵转速、阀开度是自变量，我们能调；流速是中间量；
芯片温度是因变量，是要稳的目标。}

\shuo{第二，两个待辨识的函数——函数 f 是「执行器到流速」，泵和阀怎么决定流速，我们不清楚；
函数 g 是「流速、水温、热流到芯片温度」，散热片的特性，我们也不清楚。}

\begin{yaodian}
\textbf{这两个函数不清楚，控制算法就没法设计。}\\
这就是控制实验要解决的问题。
\end{yaodian}

\subsection{进阶阅读带来的一个转向}

\dongzuo{1.5 分钟}

\shuo{我原本打算在 PID 之上直接叠加 DRL。读完三篇文献后，改成了
「先辨识、后控制、暂不上 DRL」。三个理由：}

\shuo{第一，改的是顺序。Brunton 明确指出系统辨识是数据驱动控制的前置步骤——
没有可信的对象模型，任何控制器的评价都缺乏基准。这正好解释了组长为什么要求先做控制实验：
不是走流程，而是算法无法凭空设计。}

\shuo{第二，改的是对 DRL 的定位。我原来把它当「更高级的 PID」。
读完后认识到它是需要大量交互样本的决策算法。我们目前只有 CFD 仿真环境，
训练出来的策略迁移到真实台架上有明显风险。}

\shuo{第三，改的是路线。定为一阶物理模型打底、实测数据辨识参数、经典控制先行、
视数据积累情况再评估学习型控制。这条路线更慢，但更稳。}

\subsection{关键结果}

\dongzuo{2.5 分钟，切到图 2 —— 这是重点}

\shuo{这是增量式 PID 的温控仿真时程。三幅图分别是温度、阀开度、误差与积分项。}

\shuo{先说工况：目标温度 68 度，热源是 82 加减 12 度、周期 15 秒的波动。
PID 参数用的是 Kp 等于 8、Ki 等于 0.5、Kd 等于 2，控制周期 50 毫秒。}

\shuo{再看结果：芯片温度范围 50.7 到 72.1 度，稳态均值 68.7 度——
和目标 68 度差 0.7 度，基本跟上了；阀开度跑满了 0 到 100\% 的全行程，均值 54.4\%；
峰值超调正 3.9 度，也就是 5.7\%。}

\dongzuo{指中间那幅图}

\shuo{请大家注意中间这幅图。峰值那一段，阀已经开到了 100\%，
而且持续了 29.3\% 的时间。}

\dongzuo{指上面那幅图}

\shuo{同时，上面这幅图显示——温度在这段时间里还在往上走。}

\shuo{这两件事合起来说明了一个问题，就是我下面要讲的「与预期的差异」。}

\subsection{与预期的差异 [重点]}

\dongzuo{2.5 分钟 —— 这是最能体现思考深度的一页，一定要讲透}

\shuo{我原来的预期是：增量式 PID 加入积分限幅以后，应该能消除积分饱和引起的超调。}

\shuo{实际结果不是这样。}

\shuo{限幅确实起作用了——积分项峰值从 1052.6 被压到了 100。
但超调并没有消失，峰值仍然有正 3.9 度。限幅只是把超调从「持续增大」变成了「有上界」。}

\shuo{那真正的原因是什么？}

\shuo{我把误差、积分项、阀开度三条曲线叠在一起看，才找到答案：
根本原因不是积分饱和，而是执行器的散热能力存在物理上限。}

\begin{yaodian}
\textbf{仿真里有 29.3\% 的时间，阀开度已经到 100\% 了，但误差仍然是正的。}\\
这个时候，无论控制器输出什么，都不会有效果——\\
\textbf{温度必然继续上升，形成结构性的超调。}
\end{yaodian}

\shuo{这个差异改变了我的下一步方向。我不再继续在 PID 参数和抗饱和上做文章了，
而是必须先把执行器的实际能力边界测出来。}

\shuo{举个很具体的例子：如果台架上泵的最大流量不足以应对峰值热流，
那这就是选型问题，不是算法问题。再怎么调 PID 也没用。}

\shuo{这也再次指向了组长要求的控制实验。}

\subsection{问题与对策}

\dongzuo{2 分钟}

\shuo{我列了三个主要问题，每个都有对应的尝试和对策：}

\shuo{问题一，执行器特性完全未知。模型里的 k(u) 是我凭经验设的幂次式，
泵转速和流速的真实关系也不清楚。参数不可信，所以仿真结论只能定性。
尝试过用已有的 CFD 结果反推散热能力——但 CFD 给的是稳态场，缺少动态响应信息，
只能给稳态趋势，拿不到时变特性。对策是设计并执行辨识实验。}

\shuo{问题二，未来要加水阀，双执行器的耦合关系不明。关小阀门会改变管路阻力，
可能影响泵的工作点。尝试过「固定阀门、只调泵」来回避耦合——
但这属于暂时回避，不是解决，水阀一装就得重做。对策是在加装水阀前先用单执行器
把泵的特性曲线测准；水阀到位后只需补测耦合项，不必重做全部工作。}

\shuo{问题三，温度测点与控制目标不一致。要稳的是芯片最高温，但传感器只能给离散测点。
尝试过用测点平均温度当控制目标——但平均温度会掩盖热点，和「防止芯片过热」的目标不符。
对策是把控制目标改成「测点最高温不超过阈值」，而不是「平均温度稳定在某值」。
这样即使测点和真实热点有偏差，也属于偏保守的方向，安全性更好。}

\subsection{阶段性结论}

\dongzuo{1.5 分钟}

\textbf{可以确定的有四条：}

\begin{enumerate}
  \item 控制回路的基本逻辑是正确的——温度偏高就增大开度，方向始终正确，没出现反向调节
  \item 增量式 PID 的两重限幅是必要的——积分限幅和输出限幅都按预期动作，
        避免了输出越界这类最危险的情况
  \item 智能算法现阶段不应作为主线——依据是数据需求和实测条件的\textbf{双重不匹配}，
        不是算法本身不好
  \item 必须先辨识、后控制——这一点已经被本月「超调并非源于积分饱和」的分析印证了。
        \textbf{对象和执行器不清楚的时候，调参就是在猜}
\end{enumerate}

\textbf{尚不能确定的有三条：}

\begin{enumerate}
  \item 执行器的能力边界——泵的最大流量、最小可调流量，以及未来水阀的分辨率，都要实测。
        \textbf{如果最大流量不足以带走峰值热流，就要更换选型，这会改变整个控制方案的前提}
  \item 两个执行器的耦合程度——这决定了后面用单变量控制还是多变量控制
  \item 控制通道的时间尺度——从改变泵转速到芯片温度出现可测变化，
        这个延迟是零点几秒还是几十秒，\textbf{直接决定控制周期怎么选}
\end{enumerate}

\subsection{下月计划 [重点]}

\dongzuo{2 分钟 —— 这一页是给老师看的重点，落在方案上}

\shuo{下个月的核心任务就是把这两个函数测出来。}

\shuo{关于函数 f，执行器到流速。如果把它当成一个完全未知的二元函数，
那要扫的是二维网格，比如 5 乘 5 就是 25 组工况。}

\shuo{我建议用增量叠加的形式来降实验量：}

\[
\Delta Q = a \cdot \Delta P + b \cdot \Delta V + c \cdot \Delta P \cdot \Delta V
\]

\shuo{这样只需要：固定阀、扫泵，定出系数 a，约 5 组；固定泵、扫阀，定出系数 b，约 5 组；
再取几组泵阀组合，验证耦合项 c，约 3 组。13 组就能定出执行器特性，比二维全扫省一半以上。}

\shuo{关于函数 g，流速和水温到芯片温度。固定热流密度扫流速，能得到稳态增益 kQ；
固定流速扫入口水温，得到 kT；再改变热流密度验证 kq。}

\begin{tishi}
\textbf{最后补充一点，报告里我注明了：}\\
\textbf{水阀是未来计划加装的，现有系统还没有这个执行器。}
所以本月的方案是按「泵为主执行器」设计的，预留了水阀接入后的耦合项补测环节，
避免加装以后要重做全部标定。
\end{tishi}

\section{收尾}

\dongzuo{30 秒}

\shuo{小结一下这两次报告。}

\shuo{阶段一我解决的是「该用什么方法」的问题——结论是先把控制通道搞清楚，
暂时不上智能算法。}

\shuo{阶段二我解决的是「这条通道具体是什么」的问题——把它整理成两个待辨识的函数，
并给出了实验方案。}

\begin{yaodian}
\textbf{下一步最关键的一件事，是把执行器的能力边界测出来。}\\
因为如果泵的流量不够，那后面的算法设计都建立在错误的前提上。
\end{yaodian}

\shuo{我的汇报到这里，请老师和同学批评指正。}

\newpage
\appendix
\renewcommand{\thesection}{附录 \Alph{section}}

\section{预计被问到的问题}

\subsection{为什么不用强化学习？别人都在用。}

\shuo{不是算法不好，是我们的条件不匹配。DRL 需要大量交互样本，
我们现在只有 CFD 仿真环境，样本来自仿真而不是实物。Vignon 的综述里也提到，
DRL 在流动控制领域训练成本高、样本效率低。}

\shuo{更关键的是——我们连执行器到流速的映射都还不知道。
在这种情况下上 DRL，是在一个不清楚的环境里训练，学出来的策略没有物理依据。}

\shuo{我的判断是：先把辨识做完，积累实测数据，再评估要不要上学习型控制。
这条路更慢，但更稳。}

\subsection{你这个仿真可信吗？参数是不是编的？}

\begin{tishi}
\textbf{这个问题问得很对，要如实说明，不要辩解。}
\end{tishi}

\shuo{仿真里的参数有一部分是假定的，特别是执行器的散热系数 k(u) 那个幂次式，
是我凭经验设的，没有实测依据。}

\shuo{所以我在报告里明确写了：参数不可信，仿真结论只能定性，不能定量。}

\shuo{这份仿真的作用不是给出准确的温度预测，而是验证控制逻辑和暴露结构性矛盾——
比如「超调来自执行器饱和而不是算法」这个结论，是定性成立的，不依赖具体参数。}

\shuo{真正的定量结果要等辨识实验做完。}

\subsection{为什么目标温度是 68 度？这个数怎么定的？}

\shuo{这个数是沿用之前仿真里的设定，没有特别的物理依据。}

\shuo{实际应该由芯片的结温上限倒推——比如留出安全余量以后的允许壳温。
这一条我需要跟结构组确认。}

\subsection{两个执行器为什么要用增量叠加形式？会不会漏掉非线性？}

\shuo{确实会。增量叠加本质上是一种局部线性化，适合在工作点附近的小范围调节。}

\shuo{我这么设计的原因是实验成本：完全未知的二元函数要扫二维网格，做 25 组以上；
增量叠加只要 13 组。}

\shuo{我的打算是分两步：先用增量叠加把工作点附近的特性测出来，够控制器用；
如果发现耦合项 c 显著，说明非线性强，再加密工况点。这样既省实验又不漏掉主要矛盾。}

\subsection{水阀还没装，你现在做这个方案有意义吗？}

\shuo{有意义，因为方案里已经预留了水阀的位置。}

\shuo{具体说：我按「泵为主执行器」设计，把泵的特性曲线先测准。
水阀到位以后，只需要补测耦合项 c，前面测的泵特性直接能用，不用重做。}

\shuo{如果我现在不做，等水阀装好了再从头设计，反而更慢。}

\subsection{那三个「尚不能确定」的问题，你打算怎么解决？}

\begin{enumerate}
  \item \textbf{执行器能力边界}——靠辨识实验。这也是下个月的第一件事。
  \item \textbf{泵阀耦合程度}——靠一个专项实验：固定泵转速、改变阀开度，
        观察泵的工作点是否移动。
  \item \textbf{控制通道时间尺度}——做一个阶跃响应实验：突然改变泵转速，
        记录温度的响应曲线，从曲线里读出延迟和时间常数。
\end{enumerate}

\shuo{这三个都要实验条件，所以我需要先跟老师和结构组确认：
泵的型号是什么、有没有流量计、温度采样频率能到多少。}

\subsection{你觉得这个项目最大的风险是什么？}

\begin{yaodian}
\textbf{执行器选型可能不满足需求。}
\end{yaodian}

\shuo{从仿真结果看，峰值热流那段时间阀已经开到 100\% 了，温度还在涨。
如果实际台架上泵的最大流量也不够，那就是硬件选型问题，算法层面解决不了。}

\shuo{这个风险要在做算法之前就排除掉——这也是我为什么把「测执行器能力边界」
排在第一位的原因。}

\subsection{阶段一你说复现发现公式偏保守 11\%，这个有什么用？}

\shuo{直接影响我怎么评价控制器。}

\shuo{如果我用理论公式算调节时间，会高估需要的响应时间，
可能把 PID 调得过于激进。用仿真量取值更接近实际。}

\shuo{另外这也说明理论公式有适用范围——调节时间公式是近似式，
在阻尼比小于 0.7 时不够准。用之前要知道它的边界在哪。}

\section{讲的时候的注意事项}

\subsection{时间分配（总共 22 分钟）}

\begin{table}[htbp]
\centering
\begin{tabularx}{\textwidth}{@{}Y c@{}}
\toprule
\rowcolor{cLight} 内容 & 时间 \\
\midrule
开场（两份报告的分工） & 0.5 分钟 \\
阶段一：阅读范围 + 核心概念 & 3.5 分钟 \\
阶段一：复现 + 收获疑问 & 4 分钟 \\
阶段二：任务 + 模型整理 & 3.5 分钟 \\
\rowcolor{cWarnBg} \textbf{阶段二：关键结果 + 与预期的差异 [重点]} & \textbf{5 分钟} \\
阶段二：问题对策 + 结论 & 3.5 分钟 \\
\rowcolor{cWarnBg} \textbf{阶段二：下月计划（实验方案）[重点]} & \textbf{2 分钟} \\
\bottomrule
\end{tabularx}
\end{table}

\begin{tishi}
\textbf{两个[重点]是重点，一定要留够时间。}如果时间不够，先砍阶段一的「核心概念」。
\end{tishi}

\subsection{三个提醒}

\begin{enumerate}
  \item \textbf{「与预期的差异」那页一定要讲透}——这是最能体现思考深度的地方。\\
        核心句式：\textbf{「我原以为 X，实际是 Y，原因是 Z，所以我改成做 W。」}

  \item \textbf{「未开始」的那一项要主动说。}执行器实测辨识我没做，
        要主动讲明「缺实验条件」，\textbf{不要含糊过去}。主动说比被问出来好。

  \item \textbf{仿真参数的局限性要主动交代。}「参数是不是编的」这个问题几乎一定会被问到。
        \textbf{与其被质疑，不如自己先说}——「参数有一部分是假定的，所以结论只能定性」。
\end{enumerate}

\subsection{三句话记住全篇}

\begin{yaodian}
如果要极简概括这两份报告：\\[4pt]
\textbf{阶段一：}控制效果的瓶颈不在算法，在对控制通道的了解程度。\\[2pt]
\textbf{阶段二：}把控制通道整理成两个待辨识的函数，
并发现超调来自执行器饱和而非算法。\\[2pt]
\textbf{下一步：}先测执行器能力边界——这是所有后续工作的前提。
\end{yaodian}

\vfill
\begin{center}\small\color{cGrey}
—— 分享稿 v1.0 ｜ 对应报告：阶段一\_控制算法（6 页）、阶段二\_控制算法（7 页）——
\end{center}

\end{document}
